\begin{examplebox}
	\textbf{\textcolor{CExample}{例 1（直接套用）}}
	\textbf{\textcolor{CExample}{题目：}}
	求函数 $y=\dfrac{2x+1}{{x-1}}$ 的值域。
	\textbf{\textcolor{CExample}{解题步骤：}}
	\begin{description}[style=nextline, leftmargin=2.8em, labelsep=0.5em, itemsep=0.45em]
		\item[\textbf{第一步（配凑分子）：}]
			将分子写成分母的倍数加余数：
			\[
				2x+1 = 2(x-1) + 3.
			\]
			\par\textbf{理由：} 目的是构造 $(x-1)$ 项，为分离常数做准备。
		\item[\textbf{第二步（分离常数）：}]
			代入原函数：
			\[
				\begin{aligned}
					y &= \frac{2(x-1)+3}{{x-1}} \\ &= 2 + \frac{3}{{x-1}}.
			\end{aligned}
			\]
			\par\textbf{理由：} 分式拆成常数与反比例函数之和。
		\item[\textbf{第三步（直接得值域）：}]
			由 $x-1\neq0$ 知 $\dfrac{3}{{x-1}}\neq0$，故 $y\neq2$。
			\par\textbf{理由：} 反比例函数值域为非零实数，平移常数 $2$ 后值域排除 $2$。
	\end{description}
	\textbf{\textcolor{CExample}{关键结论：}}
	值域为 $\{y\in\mathbb{R}\mid y\neq2\}$。

	\medskip
	\textbf{\textcolor{CExample}{例 2（变形化简）}}
	\textbf{\textcolor{CExample}{题目：}}
	求函数 $y=\dfrac{x+3}{{2x-4}}$ 的值域。
	\textbf{\textcolor{CExample}{解题步骤：}}
	\begin{description}[style=nextline, leftmargin=2.8em, labelsep=0.5em, itemsep=0.45em]
		\item[\textbf{第一步（化简系数）：}]
			将分母提出因子 $2$：$y=\dfrac{x+3}{{2(x-2)}}$。
			\par\textbf{理由：} 分离常数前先整理分母以简化后续计算。
		\item[\textbf{第二步（分离常数）：}]
			对 $\dfrac{x+3}{{x-2}}$ 分离常数：
			\[
				\dfrac{x+3}{{x-2}} = 1 + \dfrac{5}{{x-2}}.
			\]
			代入得
			\[
				\begin{aligned}
					y &= \frac{1}{2}\left(1 + \frac{5}{{x-2}}\right) \\ &= \frac12 + \frac{5}{{2(x-2)}}.
			\end{aligned}
			\]
			\par\textbf{理由：} 分子除以分母得商和余数。
		\item[\textbf{第三步（得值域）：}]
			由于 $\dfrac{5}{{2(x-2)}}\neq0$（$x\neq2$），故 $y\neq\dfrac12$。
			\par\textbf{理由：} 反比例部分非零，整体排除常数项。
	\end{description}
	\textbf{\textcolor{CExample}{关键结论：}}
	值域为 $\{y\in\mathbb{R}\mid y\neq\frac12\}$。

	\medskip
	\textbf{\textcolor{CExample}{例 3（参数分析）}}
	\textbf{\textcolor{CExample}{题目：}}
	已知函数 $y=\dfrac{3x+m}{{x+2}}$ 的值域为 $\{y\in\mathbb{R}\mid y\neq3\}$，求实数 $m$ 的取值范围。
	\textbf{\textcolor{CExample}{解题步骤：}}
	\begin{description}[style=nextline, leftmargin=2.8em, labelsep=0.5em, itemsep=0.45em]
		\item[\textbf{第一步（识别参数）：}]
			对照标准形式 $y=\dfrac{ax+b}{{cx+d}}$，得 $a=3,b=m,c=1,d=2$。
			\par\textbf{理由：} 便于套用值域公式。
		\item[\textbf{第二步（应用值域条件）：}]
			值域条件 $\{y\neq3\}$ 即 $y\neq a/c=3$，这自动满足前提 $ad\neq bc$ 下的结论。但还需检验是否可能为常数函数。
			\par\textbf{理由：} 值域排除 $a/c$ 是结论直接结果，但需避免 $ad=bc$ 的特殊情况。
		\item[\textbf{第三步（限制常数情形）：}]
			计算 $ad = 3\cdot2=6$，$bc = m\cdot1=m$。若 $ad=bc$，即 $m=6$，则函数可化为常数 $y=3$，值域为 $\{3\}$，与题设不符。故 $m\neq6$。
			\par\textbf{理由：} 当 $ad=bc$ 时原式为常数函数，值域为单点集。
	\end{description}
	\textbf{\textcolor{CExample}{关键结论：}}
	$m$ 的取值范围是 $m\in\mathbb{R}$ 且 $m\neq6$。
\end{examplebox}
